\honest{GAZP vs. GLUT}{Eliezer Yudkowsky}{2008}

Daniel Dennett reports seeing only one defence of zombies, and it rests on a descendant of ``Ned Block's Giant Lookup Table fantasy.'' A giant lookup table stores an answer for every possible input. One for a human brain would not fit in our universe, but philosophers may suppose it anyway.

Is such a table a zombie, something that behaves exactly like a human without being conscious? I cannot recall anyone claiming it is conscious; it is the reductio ad absurdum for anyone who says consciousness is simply an input-output pattern. \nb{The reductio is Block's own, from 1981, where Block used the machine against behaviorism about intelligence. Shannon and McCarthy had raised the idea in 1956.}

So where do its papers about consciousness come from? I answer with a game I call ``Follow-The-Improbability'': an improbable belief, or an improbable table, needs a source. If the table was computed from a specification of a human brain, it writes ``because of a conscious algorithm,'' and it is no more a zombie than a cellphone. \nb{Block made the same comparison in 1981, with a two-way radio: the machine is ``a conduit for intelligence.'' ``A conscious algorithm'' assumes that a brain's specification run elsewhere is conscious, which the previous post left for later. The argument does not need it, since the specification came from a human.} A few paragraphs later I say that ``The consciousness isn't in the lookup table.'' \nb{So in the sense of the opening question the table is still a zombie. The cellphone answer is about where its talk comes from.}

If the table was generated at random by a quantum device, the device, I say, only splits the universe into branches, one for each outcome. \nb{Yudkowsky states the many-worlds reading as fact here. A random table would be just as improbable on any reading.} Someone must still have picked the right table out of the bin, and that someone is ``probably conscious.'' If the philosopher stipulates pure chance, the improbability lies in the philosopher's specification, and in the readers who imagine what the table would say.

My moral: trace talk about consciousness back and you ``generally'' find consciousness. I admit that this is ``not standard philosophical procedure,'' which lets thought experiments ignore how likely they are. And if someone really did draw the right table by pure chance? ``Well, then it wouldn't be conscious. IMHO.'' \nb{So, on the post's view, a table drawn by pure chance would talk as we do without being conscious. The post's reply to the philosopher is that such a thing would not arise, not that it could not exist.}

In a coda I apply the same reasoning to people who think an arbitrary AI would be moral, and I say, as a joke about coincidence, that this is why the method matters to my work.
